% Part: first-order-logic
% Chapter: axiomatic-deduction
% Section: soundness

\documentclass[../../../include/open-logic-section]{subfiles}

\begin{document}

\iftag{FOL}
      {\olfileid{fol}{axd}{sou}}
      {\olfileid{pl}{axd}{sou}}

\olsection{বিশুদ্ধতা}

\begin{explain}
স্বতঃসিদ্ধমূলক নিষ্পাদনের মতো কোনো !!^a{derivation} পদ্ধতি ঠিক তখনই
\emph{বিশুদ্ধ}, যখন বাস্তবে সত্য নয় এমন কিছু সেটি !!{derive} করতে
পারে না। তাই বিশুদ্ধতা !!{derivation} পদ্ধতির এক ধরনের নিশ্চিত
নিরাপত্তা-ধর্ম। কোন প্রমাণতাত্ত্বিক ধর্মটি বিবেচ্য তার ওপর নির্ভর করে
আমরা যেমন জানতে চাই যে
\begin{enumerate}
\item প্রত্যেক !!{derivable}~$!A$ বৈধ;
\item $!A$ অন্য কয়েকটি~$\Gamma$ থেকে !!{derivable} হলে সেটি তাদের
  অনুসিদ্ধান্তও;
\item !!{formula}s-এর কোনো সেট~$\Gamma$ অসঙ্গত হলে সেটি
  অপরিতৃপ্তিযোগ্য।
\end{enumerate}
এগুলি কোনো !!a{derivation} পদ্ধতির গুরুত্বপূর্ণ ধর্ম। এর কোনোটি সত্য
না হলে !!{derivation} পদ্ধতিটি ত্রুটিপূর্ণ—সেটি অত্যধিক কিছু
!!{derive} করত। তাই কোনো !!a{derivation} পদ্ধতির বিশুদ্ধতা প্রতিষ্ঠা
করা অত্যন্ত গুরুত্বপূর্ণ।
\end{explain}

\begin{prop}
  $!A$ কোনো স্বতঃসিদ্ধ হলে
  \iftag{FOL}
    {প্রত্যেক !!{structure}~$\Struct{M}$ ও চলরাশি-আরোপ~$s$-এর জন্য $\Sat{M}{!A}[s]$।}
    {প্রত্যেক !!{valuation}~$\pAssign{v}$-এর জন্য $\pSat{v}{!A}$।}
\end{prop}

\begin{proof}
\iftag{FOL}{সব স্বতঃসিদ্ধ বৈধ কি না তা যাচাই করতে হবে। যেমন,
  \olref[qua]{ax:q1}-এর ক্ষেত্রটি দেখি: ধরা যাক $t$, $!A$-তে
  $x$-এর জন্য !!{free for}, এবং $\Sat{M}{\lforall[x][!A]}[s]$।
  তাহলে পরিতৃপ্তির সংজ্ঞা অনুসারে প্রত্যেক $\varAssign{s'}{s}{x}$-এর
  জন্য $\Sat{M}{!A}[s']$; বিশেষত $s'(x) = \Value{t}{M}[s]$ হলেও এটি
  সত্য। \olref[syn][ext]{prop:ext-formulas} অনুসারে
  $\Sat{M}{\Subst{!A}{t}{x}}[s]$। তাই
  $\Sat{M}{(\lforall[x][!A] \lif \Subst{!A}{t}{x})}[s]$।}{প্রত্যেক
  স্বতঃসিদ্ধের সত্যসারণি তৈরি করে যাচাই করো যে সেগুলি সর্বতঃসত্য।}
\end{proof}

\begin{thm}[বিশুদ্ধতা]
\ollabel{thm:soundness}
$\Gamma \Proves !A$ হলে $\Gamma \Entails !A$।
\end{thm}

\begin{proof}
$\Gamma$ থেকে~$!A$-এর !!{derivation}-এ অনুমান-বিধি দিয়ে সমর্থিত
ধাপের সংখ্যার ওপর আরোহ প্রয়োগ করি। অনুমান-বিধি দিয়ে কোনো ধাপ সমর্থিত না
হলে !!{derivation}-এর সব !!{formula}s হয় স্বতঃসিদ্ধের রূপ, নয়তো
$\Gamma$-তে আছে। আগের প্রতিজ্ঞা অনুযায়ী সব স্বতঃসিদ্ধ
\iftag{FOL}{বৈধ}{সর্বতঃসত্য}; তাই $!A$ স্বতঃসিদ্ধ হলে $\Gamma
\Entails !A$। আর $!A \in \Gamma$ হলে তুচ্ছভাবেই $\Gamma \Entails !A$।
% BN-SRC-067 TARGET-CORRECTION-BEGIN
\par\smallskip\noindent\textnormal{\small\textbf{উৎস-সংশোধন (BN-SRC-067):} হিমায়িত ইংরেজি প্রমাণটি আরোহের মাপকাঠিকে !!{derivation}-এর “দৈর্ঘ্য” বলে, কিন্তু ভিত্তিতে অনুমান-বিধি দিয়ে সমর্থিত শূন্য ধাপ এবং ধাপে এমন আগের অংশ ব্যবহার করে যাতে অন্তত একটি কম অনুমান-সমর্থিত ধাপ আছে। বাংলা পাঠে প্রমাণটি যে মাপকাঠি বাস্তবে ব্যবহার করে—অনুমান-বিধি দিয়ে সমর্থিত ধাপের সংখ্যা—সেটিই বলা হয়েছে।} % BN-SRC-067 TARGET-CORRECTION-END

~$!A$-এর !!{derivation}-এর শেষ ধাপটি মোডাস পোনেন্স দিয়ে সমর্থিত
হলে !!{derivation}-এ !!{formula}s $!B$ ও~$!B \lif !A$ আছে এবং ওই
!!{formula}s-এ শেষ হওয়া !!{derivation}-এর অংশগুলির ওপর আরোহের অনুমান
প্রযোজ্য (কারণ সেগুলিতে অনুমান দিয়ে সমর্থিত অন্তত একটি কম ধাপ আছে)।
তাই আরোহের অনুমান অনুসারে $\Gamma \Entails !B$ এবং $\Gamma \Entails
!B \lif !A$। অতএব
\iftag{FOL}
      {\olref[syn][sem]{thm:sem-deduction}}
      {\olref[pl][syn][sem]{thm:sem-deduction}}
অনুসারে $\Gamma \Entails !A$।

\iftag{FOL}{এবার ধরা যাক, শেষ ধাপটি~\QR দিয়ে সমর্থিত। তাহলে ধাপটির
রূপ $!C \lif \lforall[x][!B(x)]$ এবং তার আগে $!C \lif !B(c)$ ধাপটি
আছে, যেখানে $c$, $\Gamma$, $!C$ বা~$\lforall[x][!B(x)]$-এ নেই।
আরোহের অনুমান অনুসারে $\Gamma \Entails !C \lif !B(c)$।
\olref[syn][sem]{thm:sem-deduction} থেকে $\Gamma \cup \{!C\}
\Entails !B(c)$।
% BN-SRC-068 TARGET-CORRECTION-BEGIN
\par\smallskip\noindent\textnormal{\small\textbf{উৎস-সংশোধন (BN-SRC-068):} হিমায়িত ইংরেজি \QR-ক্ষেত্রের প্রথম বাক্যে সর্বজনীন সূত্রের দুটি লেখায় সূত্র-চলকটির উপসর্গচিহ্ন বাদ দিয়ে $\lforall[x][B(x)]$ লেখা হয়েছে, যদিও পূর্ববর্তী ধাপ ও বাকি প্রমাণে সূত্রটি $!B(x)$। বাংলা পাঠে দুই জায়গাতেই $\lforall[x][!B(x)]$ পুনরুদ্ধার করা হয়েছে।} % BN-SRC-068 TARGET-CORRECTION-END

এমন কোনো গঠন~$\Struct{M}$ বিবেচনা করো যেন $\Sat{M}{\Gamma \cup
  \{!C\}}$। দেখাতে হবে, $\Sat{M}{\lforall[x][!B(x)]}$। যেহেতু
$\lforall[x][!B(x)]$ একটি !!a{sentence}, তার মানে প্রত্যেক
চলরাশি-আরোপ~$s$-এর জন্য $\Sat{M}{!B(x)}[s]$ দেখাতে হবে
(\olref[syn][ass]{prop:sat-quant})। $\Gamma \cup \{!C\}$ শুধু বাক্য
দিয়ে গঠিত বলে \olref[syn][sat]{defn:satisfaction} অনুসারে প্রত্যেক
$!D \in \Gamma$-র জন্য $\Sat{M}{!D}[s]$। $\Struct{M'}$-কে
$\Struct{M}$-এর মতোই ধরি, শুধু $\Assign{c}{M'} = s(x)$। যেহেতু
$c$, $\Gamma$ বা~$!C$-তে আসে না,
\olref[syn][ext]{cor:extensionality-sent} অনুসারে $\Sat{M'}{\Gamma
\cup \{!C\}}$। $\Gamma \cup \{!C\} \Entails !B(c)$ বলে
$\Sat{M'}{!B(c)}$। যেহেতু $!B(c)$ একটি !!a{sentence}, তাই
\olref[syn][ass]{prop:sentence-sat-true} অনুসারে
$\Sat{M'}{!B(c)}[s]$। \olref[syn][ext]{prop:ext-formulas} অনুসারে
$\Sat{M'}{!B(x)}[s]$ যদি এবং কেবল যদি $\Sat{M'}{!B(c)}[s]$ (মনে
করো, $!B(c)$ কেবল $\Subst{!B(x)}{c}{x}$)। তাই
$\Sat{M'}{!B(x)}[s]$। $c$, $!B(x)$-এ আসে না বলে
\olref[syn][ext]{prop:extensionality} অনুসারে
$\Sat{M}{!B(x)}[s]$। কিন্তু~$s$ ছিল যেকোনো চলরাশি-আরোপ; অতএব
$\Sat{M}{\lforall[x][!B(x)]}$। তাই $\Gamma \cup \{!C\} \Entails
\lforall[x][!B(x)]$। \olref[syn][sem]{thm:sem-deduction} অনুসারে
$\Gamma \Entails !C \lif \lforall[x][!B(x)]$।
% BN-SRC-069 TARGET-CORRECTION-BEGIN
\par\smallskip\noindent\textnormal{\small\textbf{উৎস-সংশোধন (BN-SRC-069):} হিমায়িত ইংরেজি প্রমাণে $\Gamma \cup \{!C\} \Entails !B(c)$ থেকে পরের পরিতৃপ্তি-বিবৃতিতে উপসর্গচিহ্ন বাদ দিয়ে $\Sat{M'}{B(c)}$ লেখা হয়েছে। পরের বাক্য ও প্রতিস্থাপন-সমতা উভয়েই $!B(c)$ ব্যবহৃত; বাংলা পাঠে $\Sat{M'}{!B(c)}$ রাখা হয়েছে।} % BN-SRC-069 TARGET-CORRECTION-END

$!A$ যদি~\QR দিয়ে সমর্থিত হয়, কিন্তু তার রূপ $\lexists[x][!B(x)]
\lif !C$ হয়, সেই ক্ষেত্রটি অনুশীলনী হিসেবে থাকল।}{}
\end{proof}

\tagprob{FOL}
\begin{prob}
  \olref[fol][axd][sou]{thm:soundness}-এর প্রমাণ সম্পূর্ণ করো।
\end{prob}
\tagendprob

\begin{cor}
\ollabel{cor:weak-soundness}
$\Proves !A$ হলে $!A$ \iftag{FOL}{বৈধ}{একটি সর্বতঃসত্য}।
\end{cor}

\begin{cor}
\ollabel{cor:consistency-soundness}
$\Gamma$ পরিতৃপ্তিযোগ্য হলে সেটি সঙ্গত।
\end{cor}

\begin{proof}
বিপরীত-প্রতিজ্ঞা প্রমাণ করি। ধরা যাক~$\Gamma$ সঙ্গত নয়। তাহলে
$\Gamma \Proves \lfalse$; অর্থাৎ~$\Gamma$ থেকে~$\lfalse$-এর কোনো
!!a{derivation} আছে। \olref{thm:soundness} অনুযায়ী,
\iftag{FOL}{যেকোনো !!{structure}~$\Struct{M}$}{যেকোনো !!{valuation}~$\pAssign{v}$}
যা~$\Gamma$-কে পরিতৃপ্ত করে, সেটিকে~$\lfalse$-ও পরিতৃপ্ত করতে হবে।
প্রত্যেক \iftag{FOL}{!!{structure}~$\Struct{M}$}{!!{valuation}~$\pAssign{v}$}-এর
জন্য \iftag{FOL}{$\Sat/{M}{\lfalse}$}{$\pSat/{v}{\lfalse}$} বলে কোনো
\iftag{FOL}{$\Struct{M}$}{$\pAssign{v}$} $\Gamma$-কে পরিতৃপ্ত করতে
পারে না; অর্থাৎ $\Gamma$ পরিতৃপ্তিযোগ্য নয়।
\end{proof}


\end{document}
